% Biên dịch: xelatex -shell-escape main.tex (chạy 2 lần)
% Trên Overleaf: chọn compiler XeLaTeX hoặc pdfLaTeX đều được.
\documentclass[a4paper]{article}
\usepackage{iftex}
\ifXeTeX
  \usepackage{fontspec}
\else
  \usepackage{vntex}
\fi
\usepackage{pdflscape}
\usepackage{rotating}
\usepackage{a4wide,amssymb,epsfig,latexsym,multicol,array,hhline,fancyhdr}
\usepackage{caption}
\usepackage{subcaption}
\usepackage{booktabs}
\usepackage{amsmath}
\usepackage{float}
\usepackage{lastpage}
\usepackage{enumerate}
\usepackage{color}
\usepackage{graphicx}					% Standard graphics package
\usepackage{array}
\usepackage{tabularx, caption}
\usepackage{multirow}
\usepackage[framemethod=tikz]{mdframed}% For highlighting paragraph backgrounds
\usepackage{multicol}
\usepackage{rotating}
\usepackage{graphics}
\usepackage{wrapfig}
\usepackage{geometry}
\usepackage{setspace}
\usepackage{epsfig}
\usepackage{tikz}
\usepackage{listings}
\usetikzlibrary{arrows,snakes,backgrounds,positioning,arrows.meta,fit}
\usepackage{hyperref}
% minted v3 cần latexminted, bản đi kèm Ubuntu bị lỗi với Python 3.14 nên ưu tiên minted2
\IfFileExists{minted2.sty}{\usepackage{minted2}}{\usepackage{minted}}
\usepackage{indentfirst}
\usepackage{cases}
\IfFileExists{algpseudocode.sty}{\usepackage{algorithm}\usepackage{algpseudocode}}{}
\renewcommand{\baselinestretch}{1.5}
\usepackage{xcolor}
\renewcommand{\figurename}{Hình}

%%%%%%%%
\usepackage{longtable}
\usepackage[shortlabels]{enumitem}
\hypersetup{urlcolor=blue,linkcolor=black,citecolor=black,colorlinks=true}

\setlength{\headheight}{40pt}
\pagestyle{fancy}
\fancyhead{} % clear all header fields
\fancyhead[L]{
 \begin{tabular}{rl}
    \begin{picture}(25,15)(0,0)
    \put(0,-8){\includegraphics[width=8mm, height=8mm]{images/khtn.png}}
    \end{picture}&
	\begin{tabular}{l}
		\textbf{\bf \ttfamily Trường đại học Khoa Học Tự Nhiên TP.HCM}\\
		\textbf{\bf \ttfamily Khoa Công Nghệ Thông Tin}
	\end{tabular}
 \end{tabular}
}
\fancyhead[R]{
	\begin{tabular}{l}
		\tiny \bf \\
		\tiny \bf
	\end{tabular}  }
\fancyfoot{} % clear all footer fields
\fancyfoot[L]{\scriptsize \ttfamily Hệ điều hành Linux và ứng dụng}
\fancyfoot[R]{\scriptsize \ttfamily Trang {\thepage}/\pageref{page:lastpage}}
\renewcommand{\headrulewidth}{0.3pt}
\renewcommand{\footrulewidth}{0.3pt}

%%%
\setcounter{secnumdepth}{4}
\setcounter{tocdepth}{3}

%%% Khung code cho lệnh và file cấu hình
\setminted{
  frame=single,
  bgcolor=lightgray!20,
  fontsize=\footnotesize,
  baselinestretch=1,
  breaklines=true,
  breakanywhere=true,
}
\renewcommand{\listingscaption}{File}
\renewcommand{\listoflistingscaption}{Danh sách file}

% \confile{ngôn ngữ}{file trong configs/}{đường dẫn trên server}{nhãn}
\newcommand{\confile}[4]{%
  \begin{listing}[H]
    \inputminted{#1}{configs/#2}
    \caption{\texttt{#3}}
    \label{#4}
  \end{listing}}

% \shot{file}{độ rộng theo \textwidth}{chú thích}{nhãn}
\newcommand{\shot}[4]{%
  \begin{figure}[H]
    \centering
    \fbox{\includegraphics[width=#2\textwidth]{images/#1}}
    \caption{#3}
    \label{#4}
  \end{figure}}

\newcommand{\cmd}[1]{\texttt{#1}}
\setlength{\emergencystretch}{3em} % cho phép ngắt dòng quanh các lệnh dài
\definecolor{okgreen}{rgb}{0.1,0.5,0.2}
\newcommand{\ok}{\textcolor{okgreen}{\textbf{Đạt}}}

\begin{document}

\begin{titlepage}
\begin{center}
ĐẠI HỌC QUỐC GIA THÀNH PHỐ HỒ CHÍ MINH\\
TRƯỜNG ĐẠI HỌC KHOA HỌC TỰ NHIÊN\\
KHOA CÔNG NGHỆ THÔNG TIN
\end{center}

\vspace{0.7cm}

\begin{figure}[h!]
\begin{center}
\includegraphics[width=5cm]{images/khtn.png}
\end{center}
\end{figure}

\vspace{0.5cm}


\begin{center}
\begin{tabular}{c}
\multicolumn{1}{c}{\textbf{{\normalsize HỆ ĐIỀU HÀNH LINUX VÀ ỨNG DỤNG}}}\\
~~\\
\hline
\\
\multicolumn{1}{l}{\textbf{{\Large Bài thực hành}}}\\
\\
\textbf{{\Large TRIỂN KHAI WEB SERVER APACHE TRÊN LINUX}}\\
\\
\hline
\end{tabular}
\end{center}

\vspace{0.5cm}

\begin{table}[h]
\begin{tabular}{rrl}
\hspace{4.5cm} & Giảng viên hướng dẫn: & Lê Hà Minh\\
\hspace{4.5cm} & Sinh viên thực hiện: & Nguyễn Thành Dự - 24120288\\

\end{tabular}
\end{table}

\vspace{2.5cm}

\begin{center} {\footnotesize Thành phố Hồ Chí Minh, 10/2026}
\end{center}
\end{titlepage}


\thispagestyle{empty}

\newpage
\renewcommand{\contentsname}{Mục lục}
\renewcommand{\tablename}{Bảng}
\renewcommand{\figurename}{Hình}
\renewcommand{\refname}{Tài liệu tham khảo}

\tableofcontents

\newpage

\renewcommand{\arraystretch}{1.3} % Adjust row height
\setcounter{page}{1}

\section{Môi trường thực hiện}

\subsection{Thông tin hệ thống}
Bài thực hành được thực hiện trên một máy ảo \textbf{AWS EC2} (vùng
\texttt{ap-southeast-1}, Singapore) chạy \textbf{Ubuntu}. Các thành phần chính:

\begin{table}[H]
  \centering
  \small
  \begin{tabularx}{\textwidth}{>{\raggedright}p{4.2cm}X}
    \toprule
    \textbf{Thành phần} & \textbf{Sử dụng} \\
    \midrule
    Web server & \cmd{apache2}, chạy dưới user \cmd{www-data} \\
    Cấu hình virtual host & \cmd{/etc/apache2/sites-available/*.conf}, bật bằng \cmd{a2ensite} \\
    Kiểm tra cấu hình & \cmd{apache2ctl configtest}, \cmd{apache2ctl -S} \\
    Module & \cmd{ssl proxy proxy\_http proxy\_balancer lbmethod\_byrequests headers} (bật bằng \cmd{a2enmod}) \\
    Chứng chỉ TLS & \cmd{/etc/ssl/certs}, \cmd{/etc/ssl/private} \\
    Backend & Node.js, chạy bằng systemd \\
    Tường lửa & ufw (trong OS) và AWS Security Group (ngoài máy ảo) \\
    Kiểm soát truy cập bắt buộc & AppArmor \\
    \bottomrule
  \end{tabularx}
  \caption{Môi trường triển khai}
\end{table}

\subsection{Sơ đồ hệ thống}
Trên AWS, mỗi request từ trình duyệt phải đi qua \textbf{hai lớp tường lửa}:
Security Group của AWS (ngoài máy ảo) rồi mới tới ufw (trong hệ điều hành).
Điểm này là nguyên nhân của sự cố ở Mục~\ref{sec:loi}.

\begin{figure}[H]
\centering
\resizebox{\textwidth}{!}{%
\begin{tikzpicture}[
    font=\small,
    box/.style={draw, rounded corners=2pt, align=center, minimum height=0.9cm, fill=white},
    fw/.style={box, fill=orange!15, minimum width=1.9cm, minimum height=2.6cm},
    be/.style={box, fill=blue!8, minimum width=2.6cm, minimum height=0.6cm},
    arr/.style={-{Stealth[length=2mm]}, thick},
    every node/.append style={execute at begin node=\setlength{\baselineskip}{1.1em}},
  ]
  \node[box, minimum width=2cm, minimum height=2.6cm] (client) {Trình duyệt\\(máy cá nhân)};
  \node[fw, right=0.5cm of client] (sg) {AWS\\Security\\Group\\[2pt]\scriptsize 22, 80, 443};
  \node[fw, right=0.5cm of sg] (ufw) {ufw\\[2pt]\scriptsize OpenSSH\\\scriptsize Apache Full};
  \node[box, right=0.5cm of ufw, minimum width=3.6cm, minimum height=2.6cm, fill=green!8] (apache)
      {\textbf{Apache2}\\[3pt]:80 $\xrightarrow{301}$ :443 (TLS)\\[3pt]
       \scriptsize vhost \texttt{demo} $\to$ /var/www/demo\\
       \scriptsize vhost \texttt{demo2} $\to$ /var/www/demo2};
  \node[be, right=0.8cm of apache, yshift=0.9cm] (n0) {\scriptsize node 127.0.0.1:3000};
  \node[be, below=0.25cm of n0] (n1) {\scriptsize node 127.0.0.1:3001};
  \node[be, below=0.05cm of n1] (n2) {\scriptsize node 127.0.0.1:3002};
  \node[be, below=0.05cm of n2] (n3) {\scriptsize node 127.0.0.1:3003};

  \draw[arr] (client) -- (sg);
  \draw[arr] (sg) -- (ufw);
  \draw[arr] (ufw) -- (apache);
  \draw[arr] (apache.east) ++(0,0.9) -- node[above, font=\scriptsize] {/api} (n0.west);
  \draw[arr] (apache.east) ++(0,-0.3) -- node[above, font=\scriptsize] {/app} (n2.west);

  \begin{scope}[on background layer]
    \node[draw, dashed, rounded corners, inner sep=6pt, fit=(ufw)(apache)(n0)(n3),
          label={[font=\scriptsize]above right:EC2 Ubuntu (ip-172-31-44-12)}] {};
    \node[draw, dashed, gray, rounded corners, inner sep=4pt, fit=(n1)(n3),
          label={[font=\scriptsize, gray]below:balancer://appcluster}] {};
  \end{scope}
\end{tikzpicture}}
\caption{Kiến trúc triển khai và hai lớp tường lửa trên AWS}
\label{fig:arch}
\end{figure}

\section{Các phần thực hiện}

\subsection{Phần 0 -- Chuẩn bị môi trường}
\textbf{\large Yêu cầu:} Cho hai tên miền \cmd{demo.hcmus.edu.vn} và
\cmd{demo2.hcmus.edu.vn} phân giải về máy chủ (lab không có DNS thật).

\textbf{\large Thực hiện:} Thêm vào \cmd{/etc/hosts} trên máy cá nhân (IP public
của EC2), và trên chính EC2 trỏ hai tên miền về localhost:
\begin{minted}{text}
# máy cá nhân
13.212.177.185  demo.hcmus.edu.vn demo2.hcmus.edu.vn
# trên EC2
127.0.0.1       demo.hcmus.edu.vn demo2.hcmus.edu.vn
\end{minted}

\subsection{Phần A -- Cài đặt và baseline}
\textbf{\large Yêu cầu:} Cài Apache, mở tường lửa, kiểm tra dịch vụ đang chạy.

\textbf{\large Thực hiện:}
\begin{minted}{bash}
apt update && apt install -y apache2
a2enmod ssl proxy proxy_http proxy_balancer lbmethod_byrequests headers
systemctl enable --now apache2
ufw allow OpenSSH && ufw allow 'Apache Full' && ufw enable
\end{minted}
Ngoài ufw, trên AWS còn phải mở inbound cổng 80 và 443 trong Security Group
của instance (Hình~\ref{fig:sg}).

\textbf{\large Kết quả:}
\shot{ss-121541.png}{0.95}{Apache đang chạy và lắng nghe cổng 80, 443}{fig:active}
\shot{ss-121603.png}{0.38}{ufw cho phép OpenSSH và Apache Full}{fig:ufw}
\shot{ss-121634.png}{1.0}{Inbound rules của Security Group trên AWS: 22, 80, 443}{fig:sg}

\subsection{Phần B -- Website đầu tiên trên Virtual Host (demo)}
\textbf{\large Yêu cầu:} Tạo website tĩnh \cmd{/var/www/demo} phục vụ cho
tên miền \cmd{demo.hcmus.edu.vn}.

\textbf{\large Thực hiện:}
\begin{minted}{bash}
mkdir -p /var/www/demo
# tạo /var/www/demo/index.html với tiêu đề "Website DEMO 1"
chown -R root:www-data /var/www/demo
chmod -R 750 /var/www/demo
a2ensite demo
apache2ctl configtest && systemctl reload apache2
\end{minted}
Nội dung đầy đủ của \cmd{demo.conf} (sau khi đã thêm HTTPS, \cmd{/api},
\cmd{/app}) ở File~\ref{lst:demo}.

\textbf{\large Kết quả:}
\shot{ss-121719.png}{0.85}{Vhost 1: \texttt{demo.hcmus.edu.vn}}{fig:demo1}

\subsection{Phần C -- Virtual Host thứ hai (demo2)}
\textbf{\large Yêu cầu:} Thêm website thứ hai trên cùng IP, Apache phân biệt
hai site bằng header \texttt{Host:}.

\textbf{\large Thực hiện:} Tạo \cmd{/var/www/demo2} tương tự Phần B, tạo
\cmd{demo2.conf} rồi \cmd{a2ensite demo2}. Site \cmd{demo2} cũng được bổ sung
một VirtualHost \cmd{:443} với chứng chỉ riêng.

\confile{apacheconf}{demo2.conf}{/etc/apache2/sites-available/demo2.conf}{lst:demo2}

\textbf{\large Kết quả:} Hai tên miền cho ra hai trang khác nhau dù cùng một IP và
một tiến trình Apache. \cmd{apache2ctl -S} liệt kê hai namevhost trên cổng 80.
\shot{ss-122236.png}{0.85}{Vhost 2: \texttt{demo2.hcmus.edu.vn}}{fig:demo2}
\shot{ss-121742.png}{0.85}{\texttt{apache2ctl -S}: hai namevhost trên cổng 80, demo trên 443}{fig:vhosts}

\subsection{Phần D -- Node API backend và Reverse Proxy (/api)}
\textbf{\large Yêu cầu:} Dựng backend Node.js chỉ lắng nghe nội bộ, Apache làm
reverse proxy cho đường dẫn \cmd{/api}.

\textbf{\large Thực hiện:} Backend \cmd{/opt/api/app.js} trả JSON cho biết
instance nào xử lý request. Backend được chạy bằng systemd dạng instantiated
service, mỗi cổng là một instance \cmd{nodeapi@<port>}.

\confile{javascript}{app.js}{/opt/api/app.js}{lst:app}
\confile{ini}{nodeapi@.service}{/etc/systemd/system/nodeapi@.service}{lst:service}

\begin{minted}{bash}
systemctl daemon-reload
systemctl enable --now nodeapi@3000
\end{minted}
Trong \cmd{demo.conf} thêm \cmd{ProxyPreserveHost}, \cmd{ProxyPass} và
\cmd{ProxyPassReverse} cho \cmd{/api} (File~\ref{lst:demo}).

\textbf{\large Kết quả:} \cmd{/api/hello} trả JSON từ \cmd{node-3000}, trong khi
truy cập thẳng cổng 3000 từ bên ngoài thì không được: mọi request buộc phải đi
qua Apache.
\shot{ss-122327.png}{0.85}{\texttt{/api/hello} trả JSON từ \texttt{node-3000}}{fig:api}
\shot{ss-122431.png}{0.42}{Từ máy ngoài không truy cập thẳng được cổng 3000}{fig:direct}

\subsection{Phần E -- Load Balancer (pool 3 backend + failover)}
\textbf{\large Yêu cầu:} Chạy 3 instance backend, cấu hình Apache chia tải
round-robin và tự né instance bị lỗi.

\textbf{\large Thực hiện:}
\begin{minted}{bash}
systemctl enable --now nodeapi@3001 nodeapi@3002 nodeapi@3003
\end{minted}
Trong \cmd{demo.conf} khai báo \cmd{balancer://appcluster} gồm 3 thành viên với
\cmd{lbmethod=byrequests} và ánh xạ đường dẫn \cmd{/app} (File~\ref{lst:demo}).
Thử failover bằng cách dừng \cmd{nodeapi@3002} rồi gửi lại request.

\textbf{\large Kết quả:} Response xoay vòng qua 3001, 3002, 3003; khi dừng 3002,
traffic tự chuyển sang 3001 và 3003 mà client không thấy lỗi.
\shot{ss-122635.png}{0.6}{Bốn instance Node chỉ lắng nghe trên 127.0.0.1}{fig:ports}
\shot{ss-114911.png}{1.0}{Load balancer xoay vòng 3 instance; khi dừng \texttt{nodeapi@3002} traffic tự chuyển sang 3001, 3003}{fig:lb}

\subsection{Phần F -- SSL/TLS (HTTPS) + redirect + HSTS}
\textbf{\large Yêu cầu:} Bật HTTPS bằng chứng chỉ tự ký, ép redirect 80 sang
443 và bật HSTS.

\textbf{\large Thực hiện:}
\begin{minted}{bash}
openssl req -x509 -nodes -newkey rsa:2048 -days 365 \
  -keyout /etc/ssl/private/demo.key -out /etc/ssl/certs/demo.crt \
  -subj "/C=VN/ST=HCM/L=HCMC/O=HCMUS/CN=demo.hcmus.edu.vn"
\end{minted}
VirtualHost \cmd{:80} chỉ còn nhiệm vụ redirect 301 sang HTTPS; VirtualHost
\cmd{:443} phục vụ nội dung tĩnh, \cmd{/api}, \cmd{/app} và gửi header HSTS.
Đây là nội dung cuối cùng của \cmd{demo.conf}:

\confile{apacheconf}{demo.conf}{/etc/apache2/sites-available/demo.conf}{lst:demo}

\textbf{\large Kết quả:} HTTP trả \cmd{301 Moved Permanently} về \cmd{https://};
HTTPS trả \cmd{200 OK} kèm \cmd{Strict-Transport-Security}.
\shot{ss-122751.png}{0.48}{HTTP bị redirect 301 sang HTTPS; HTTPS trả 200 kèm \texttt{Strict-Transport-Security}}{fig:https}

\subsection{Phần G -- Tường lửa và AppArmor: rà soát và làm cứng}
\textbf{\large Yêu cầu:} Chỉ mở đúng cổng cần thiết, giữ cơ chế kiểm soát truy
cập bắt buộc hoạt động, ẩn version Apache và thêm header bảo mật.

\textbf{\large Thực hiện:} Tường lửa chỉ mở 22, 80, 443 (Hình~\ref{fig:ufw},
\ref{fig:sg}). AppArmor giữ nguyên chế độ enforce. Header bảo mật đặt trong
\cmd{hardening.conf} (bật bằng \cmd{a2enconf hardening}); \cmd{ServerTokens},
\cmd{ServerSignature}, \cmd{TraceEnable} đặt trong \cmd{security.conf}.

\confile{apacheconf}{hardening.conf}{/etc/apache2/conf-available/hardening.conf}{lst:hardening}
\confile{apacheconf}{security.conf}{/etc/apache2/conf-available/security.conf (các dòng có hiệu lực)}{lst:security}

\textbf{\large Kết quả:} Header \texttt{Server} chỉ còn \texttt{Apache} (đã ẩn
version), ba header bảo mật xuất hiện.
\shot{ss-122833.png}{0.8}{AppArmor đang hoạt động; header \texttt{Server} đã ẩn version và có 3 header bảo mật}{fig:harden}

\section{Bảng kiểm tra tổng}

\begin{table}[H]
  \centering
  \small
  \renewcommand{\arraystretch}{1.2}
  \begin{tabularx}{\textwidth}{c>{\raggedright}p{2.4cm}>{\raggedright}X>{\raggedright}p{3.2cm}>{\centering}p{1.3cm}c}
    \toprule
    \textbf{\#} & \textbf{Hạng mục} & \textbf{Lệnh kiểm chứng} & \textbf{Kết quả} & \textbf{Hình} & \\
    \midrule
    1  & Apache chạy      & \cmd{systemctl is-active apache2} & active, LISTEN :80/:443 & \ref{fig:active} & \ok \\
    2  & Firewall mở      & \cmd{ufw status} + Security Group & Apache Full; SG 22/80/443 & \ref{fig:ufw}, \ref{fig:sg} & \ok \\
    3  & Vhost 1          & trình duyệt \cmd{demo.hcmus.edu.vn} & trang DEMO 1 & \ref{fig:demo1} & \ok \\
    4  & Vhost 2          & trình duyệt \cmd{demo2...}; \cmd{apache2ctl -S} & trang DEMO 2; 2 namevhost & \ref{fig:demo2}, \ref{fig:vhosts} & \ok \\
    5  & Reverse proxy    & \cmd{/api/hello} & JSON node-3000 & \ref{fig:api} & \ok \\
    6  & Load balancer    & lặp \cmd{curl .../app/} & xoay vòng 3001/2/3 & \ref{fig:lb} & \ok \\
    7  & Failover         & dừng \cmd{nodeapi@3002}, lặp curl & chỉ còn 3001, 3003, không lỗi & \ref{fig:lb} & \ok \\
    8  & Redirect HTTPS   & \cmd{curl -I http://demo...} & 301 $\to$ https, có HSTS & \ref{fig:https} & \ok \\
    9  & AppArmor         & \cmd{aa-status} & module loaded, enforce & \ref{fig:harden} & \ok \\
    10 & Hardening        & \cmd{curl -kI https://...} & \texttt{Server: Apache} + 3 header & \ref{fig:harden} & \ok \\
    \bottomrule
  \end{tabularx}
  \caption{Kết quả kiểm tra tổng}
\end{table}

\section{Sự cố gặp phải và cách xử lý}
\label{sec:loi}

\subsection{Không truy cập được HTTPS do Security Group thiếu cổng 443}

\textbf{\large Triệu chứng:} Sau khi bật HTTPS (Phần F), mở
\cmd{http://demo.hcmus.edu.vn} trên trình duyệt ở máy cá nhân thì bị chuyển
sang \cmd{https://} rồi \textbf{treo cho đến khi hết thời gian chờ} (connection
timed out), không hiện trang. Trong khi đó, trên chính EC2,
\cmd{curl -kI https://demo.hcmus.edu.vn/} vẫn trả \cmd{HTTP/1.1 200 OK}.

\textbf{\large Phân tích, khoanh vùng:} Kiểm tra lần lượt từ trong ra ngoài theo
đường đi của gói tin (Hình~\ref{fig:arch}):
\begin{enumerate}
  \item \textbf{Apache:} \cmd{apache2ctl configtest} báo \texttt{Syntax OK};
        \cmd{ss -tlnp | grep :443} cho thấy \cmd{apache2} đang LISTEN trên
        \cmd{*:443}; \cmd{curl} nội bộ trả 200 $\Rightarrow$ Apache và SSL hoạt động đúng.
  \item \textbf{Tường lửa trong OS:} \cmd{ufw status} có rule
        \texttt{Apache Full} (= 80 + 443) $\Rightarrow$ ufw không chặn.
  \item \textbf{Nhật ký Apache:} \cmd{tail -f /var/log/apache2/demo\_access.log}
        trong lúc tải lại trang: chỉ ghi nhận request tới cổng 80 (bị
        redirect 301), \textbf{không có} request HTTPS nào từ máy cá nhân tới
        được Apache; \cmd{demo\_error.log} cũng không có lỗi SSL
        $\Rightarrow$ gói tin bị chặn \emph{trước khi} tới máy ảo.
  \item \textbf{Kiểu lỗi:} \emph{timeout} (gói tin bị bỏ đi âm thầm) chứ không
        phải \emph{connection refused} (máy trả RST vì không có tiến trình
        nghe). Đây là dấu hiệu điển hình của một tường lửa ở giữa.
\end{enumerate}

\textbf{\large Nguyên nhân:} Trên AWS, lưu lượng vào EC2 bị lọc bởi \textbf{hai
lớp độc lập}: Security Group (ở hạ tầng AWS) và ufw (trong Ubuntu). Security
Group lúc tạo instance chỉ có inbound cho cổng 22 và 80. Mở cổng 443 trên ufw
là chưa đủ vì gói tin HTTPS đã bị Security Group loại bỏ trước khi đến được
card mạng của máy ảo. Cổng 80 vẫn mở nên trình duyệt nhận được 301, nhưng bước
tiếp theo sang 443 thì bị chặn, khiến lỗi trông như ``trang HTTPS bị treo''.

\textbf{\large Cách sửa:} Vào \textit{EC2 $\to$ Instances $\to$ Security $\to$
Security groups (launch-wizard-9) $\to$ Edit inbound rules}, thêm rule
\textbf{HTTPS, TCP 443, nguồn 0.0.0.0/0}. Không cần khởi động lại Apache hay máy
ảo; tải lại trang thì truy cập được ngay (trình duyệt cảnh báo chứng chỉ tự ký,
đây là điều bình thường trong lab). Kết quả ở Hình~\ref{fig:sg}.

\textbf{\large Bài học:}
\begin{itemize}
  \item Trên cloud, ``mở firewall'' nghĩa là mở \textbf{cả hai lớp}: firewall
        trong OS và firewall của nhà cung cấp (Security Group).
  \item \cmd{curl} trên localhost chạy được không chứng minh được dịch vụ truy
        cập được từ ngoài; cần kiểm thử từ máy khác.
  \item Phân biệt \emph{timeout} và \emph{connection refused} giúp khoanh vùng
        nhanh: timeout $\to$ firewall chặn; refused $\to$ dịch vụ không lắng nghe.
  \item Access log không có dòng nào cho request là bằng chứng rõ ràng request
        không tới được Apache.
\end{itemize}

\subsection{Các điểm phát hiện thêm khi rà soát cấu hình}

\textbf{\large Dư dấu \texttt{/} khi proxy:} JSON trả về ở Hình~\ref{fig:api} có
\texttt{"path": "//hello"}. Nguyên nhân là
\cmd{ProxyPass /api http://127.0.0.1:3000/}: vế trái không có \texttt{/} cuối
còn vế phải có, nên \cmd{/api/hello} được ghép thành
\cmd{http://127.0.0.1:3000//hello}. Cách sửa là để hai vế đồng nhất:
\begin{minted}{apacheconf}
ProxyPass        /api/ http://127.0.0.1:3000/
ProxyPassReverse /api/ http://127.0.0.1:3000/
ProxyPass        /app/ balancer://appcluster/
ProxyPassReverse /app/ balancer://appcluster/
\end{minted}

\textbf{\large Lỗi chính tả trong unit systemd:} Bản \cmd{nodeapi@.service} lúc
làm lab ghi \cmd{WantedBy=mutil-user.target} thay vì \cmd{multi-user.target}.
Các service vẫn chạy vì được khởi động thủ công bằng
\cmd{systemctl enable --now}, nhưng target sai tên nên khi EC2 khởi động lại,
backend sẽ không tự chạy và \cmd{/api}, \cmd{/app} trả 503. Cách sửa: đổi lại
thành \cmd{multi-user.target} (bản nộp kèm ở File~\ref{lst:service} đã sửa),
chạy \cmd{systemctl daemon-reload} rồi
\cmd{systemctl reenable nodeapi@300\{0,1,2,3\}}.

\section{Kết luận}

Đã triển khai thành công một web server Apache trên Ubuntu (AWS EC2) đáp ứng
đủ 10 hạng mục kiểm tra: hai virtual host trên cùng một IP, reverse proxy tới
backend Node.js chỉ lắng nghe nội bộ, load balancer ba instance có failover,
HTTPS với redirect 301 và HSTS, tường lửa chỉ mở đúng cổng cần thiết và các
header làm cứng. Bài học lớn nhất là trên môi trường cloud, tường lửa có
nhiều lớp và cần kiểm thử từ bên ngoài chứ không chỉ từ localhost.

\label{page:lastpage}


\end{document}
